% Part: incompleteness
% Chapter: incompleteness-provability
% Section: introduction

\documentclass[../../../include/open-logic-section]{subfiles}

\begin{document}

\olfileid{inc}{inp}{int}

\olsection{Pambuka}

Miturut Hilbert, sistem aksioma kanggo sawijining struktur matematika,
kayata bilangan asli, durung cukup yen ora bisa nurunake kabeh pratelan
kang bener ngenani struktur mau. Yen digandhengake karo kawigatene
sabanjure marang sistem deduksi formal, iki nuduhake yen dheweke nganggep
kita kudu njamin, upamane, yen sistem formal kang digunakake kanggo
nalar ngenani bilangan asli ora mung konsisten nanging uga
\emph{jangkep}, yaiku saben pratelan ing basane bisa
!!{derivable} utawa negasine bisa. Teorema ora jangkep kapisan saka
G\"odel nuduhake yen sistem aksioma kaya mangkono ora ana: ora ana
sistem formal kanggo aritmetika kang jangkep, konsisten, lan
!!{axiomatizable}. Malah, ora ana teori matematika kang
``cukup kuwat,'' konsisten, lan !!{axiomatizable} kang jangkep.

Tujuan Hilbert kang luwih wigati, dadi punjere programe kanggo
menehi landhesan marang matematika modern (``klasik''), yaiku nemokake bukti
kekonsistenan finitari kanggo sistem formal kang nggambarake panalaran
klasik. Mula, tumrap program Hilbert, teorema ora jangkep kapindho saka
G\"odel dadi pukulan kang luwih abot. Teorema ora jangkep kapindho bisa
diandharake kanthi umum kaya teorema kapisan. Sacara kasar, teorema iki
nyatakake yen ora ana teori aritmetika kang cukup kuwat kang bisa
mbuktekake kekonsistenane dhewe. ``Cukup kuwat'' ing kene kudu
ngemot luwih saka~$\Th{Q}$.

Gagasan kang ndhasari bukti asli teorema ora jangkep saka G\"odel
bisa ditemokake ing paradoks Epimenides. Epimenides, wong Kreta,
nyatakake yen kabeh wong Kreta tukang ngapusi; wujud paradoks kang
luwih langsung yaiku pratelan ``ukara iki salah.'' Pokoke, kanthi
ngganti bebener nganggo !!{derivability}, G\"odel bisa mbangun
!!a{sentence} formal kang, kanthi cara ora langsung, nyatakake yen
ukara iku dhewe ora !!{derivable}. Yen !!{sentence} mau !!{derivable},
teorine bakal ora konsisten. G\"odel nuduhake yen negasi saka
!!{sentence} mau uga ora !!{derivable} saka sistem aksioma kang
dirembug. (Kanggo bagean kapindho iki, G\"odel kudu nganggep yen
teori~$\Th{T}$ nduweni sipat kang diarani ``$\omega$-konsisten.''
Kekonsistenan-$\omega$ sesambungan karo kekonsistenan, nanging sipate
luwih kuwat.\footnote{Tegese, saben teori kang $\omega$-konsisten iku
konsisten, nanging kosok baline ora mesthi.} Sawetara taun sawise
G\"odel, Rosser nuduhake yen nganggep kekonsistenan lumrah saka~$\Th{T}$
wis cukup.)

Tantangan kapisan yaiku mangerteni carane mbangun !!a{sentence} kang
nyebut awake dhewe. Kanggo saben !!{formula}~$!A$ ing basa
saka~$\Th{Q}$, tetepna~$\gn{!A}$ minangka numeral kang cocog
karo~$\Gn{!A}$. Gatekna teges iki: $!A$~iku !!a{formula} ing
basa saka~$\Th{Q}$, $\Gn{!A}$~iku bilangan asli, lan $\gn{!A}$~iku
\emph{suku} ing basa saka~$\Th{Q}$. Dadi saben !!{formula}~$!A$
ing basa saka~$\Th{Q}$ nduweni \emph{jeneng},~$\gn{!A}$, yaiku
suku ing basa saka~$\Th{Q}$; iki menehi rerangka gagasan kang ngidini
!!{formula}s ing basa saka~$\Th{Q}$ ``ngomong'' ngenani
!!{formula}s liyane. Lema sabanjure diarani lema titik tetep.

\begin{lem}
Upamana $\Th{T}$ teori sembarang kang ngembangake~$\Th{Q}$, lan
$!B(x)$ iku !!{formula} sembarang kang mung variabel~$x$ bebas.
Mula ana !!a{sentence}~$!A$ saengga
$\Th{T} \Proves!A \liff !B(\gn{!A})$.
\end{lem}

Lema iki nyatakake yen kanggo saben sipat $!B(x)$, ana
!!a{sentence}~$!A$ kang nyatakake ``$!B(x)$ bener tumrap aku,''
lan $\Th{T}$ ``ngerti'' perkara iki.

Kepiye carane mbangun !!a{sentence} kaya mangkono? Gatekna versi
paradoks Epimenides saka Quine iki:
\begin{quote}
``Ngasilake pratelan salah yen didhisiki petikane'' ngasilake pratelan salah
yen didhisiki petikane.
\end{quote}
Ukara iki ora langsung nyebut awake dhewe. Ukara iki mung nyatakake
perkara ngenani objek sintaksis ing antarane tandha petik, lan kanthi
cara kuwi padha karo ukara-ukara kaya
\begin{enumerate}
\item ``Robert'' iku jeneng kang apik.
\item ``Aku mlayu.'' iku ukara cekak.
\item ``Nduweni telung tembung'' nduweni telung tembung.
\end{enumerate}
Nanging apa kang kelakon yen njupuk frasa ``ngasilake pratelan salah yen
didhisiki petikane,'' banjur frasa mau didhisiki versi kang dipetik
saka frasa iku dhewe? Banjur kita entuk ukara asal mau!{} Cekake,
ukara iki nyatakake yen ukara iku dhewe salah.

\end{document}
